% vim: set tw=80 aw sw=2 sts=2 noet:
\documentclass{beamer}

%\includeonlyframes{c} % speeding up compilation speed during debug


\usepackage[utf8x]{inputenc}            % diacritice
\usepackage[english]{babel}
\usepackage{color}                      % highlight
\usepackage{alltt}                      % highlight}        % \url{http://...} | \href{http://...}{Nume Link}

% Pentru a include cod decomentati urmatoarele 3 linii
%\usepackage{color}			 % highlight
%\usepackage{alltt}			 % highlight
%\usepackage{code/highlight}	 % highlight

\mode<presentation>
%\usetheme{UEA}
%\usepackage{imperial-beamer}
\usetheme{SCS}

% Pentru a afisa (cont.) la slide-uri prea lungi split-uite pe mai multe pag
\setbeamertemplate{frametitle continuation}[from second]
% Pentru a modifica modul de afisare al numerelor de slide
\setbeamertemplate{footline}[frame number]

\title[]{Formation Flight for Unmanned Aerial Vehicles}
\subtitle{Bachelor Thesis: July 13\textsuperscript{th} 2013}
\institute[CS]{Prof. Adina Magda Florea \\ Ing. Mihai Trăscău \\ Ș.l. Cătălin Leordeanu}
\author[A]{Alexandru George Burghelea \\ (\textit{aburghelea@rosedu.org})}

\begin{document}

{
  % Schimbam fundalul aici pentru a avea slide-ul cu logo-urile
  % Nu știu momentan cum se face asta în template.
  \usebackgroundtemplate{\includegraphics[width=\paperwidth]{title}}
  \frame{\titlepage}
}

\begin{frame}{Index}
  \begin{itemize}
    \item Intro
    \item First Chapter
    \item Second Chapter
    \item Third Chapter
    \item Results
    \item Questions
  \end{itemize}
\end{frame}

\end{document}
